\section{总结与延伸}

本讲展示了 verifier 的演进路径。Outcome verifier 通过最终答案标签支持 best-of-$N$；process reward model 用步骤监督改善信用分配和搜索；Math-Shepherd 通过 rollout 成功率自动估计步骤价值；Weaver 则组合多个弱 verifier，利用误差互补缩小 generation--verification gap，并通过蒸馏降低成本。

四种方法分别回答了四个问题：

\begin{enumerate}
    \item \textbf{结果对不对？} ORM；
    \item \textbf{从哪一步开始错？} PRM；
    \item \textbf{步骤标签能否自动产生？} rollout-based process supervision；
    \item \textbf{多个不可靠 judge 如何组合？} weak supervision ensemble。
\end{enumerate}

下一讲把反馈从“判断答案”推进到“与工具和环境交互”：ReAct 使用观察结果修正推理，RLEF 用代码执行作为训练 reward，Constitutional AI 用原则驱动模型自我批评和 AI feedback。

\end{document}
